\usepackage{makecell}

\usepackage{amsthm}
% see https://github.com/latex3/tagging-project/issues/942 and issue 920
%\ExplSyntaxOn
%\makeatletter
%\def\measuring@true{\let\ifmeasuring@\iftrue\tag_suspend:n{\measuring}\luamml_ignore:}
%\makeatother
%\ExplSyntaxOff

% \newline needs https://github.com/latex3/tagging-project/issues/1201
\newtheoremstyle{apexExample}% name
  {0pt}% Space above, empty = `usual value'
  {0pt}% Space below
  {}% Body font
  {}% Indent amount (empty = no indent, \parindent = para indent)
  {\bfseries}% Thm head font
  {}% Punctuation after thm head
  {\newline}% Space after thm head
%  {{\bfseries\thmname{#1}~\thmnumber{#2}\thmnote{#3}}}
%  {\thmname{{\bfseries#1}}\thmnumber{{ \bfseries#2}}\thmnote{{\raggedright\bfseries#3}}}
%  {\makebox[\ifbool{latexml}{10em}{8em}][l]{\thmname{#1}\thmnumber{ #2}}%
%   \thmnote{\parbox[t]{.75\textwidth}{\raggedright#3}}%
%  }
  {\thmname{\makebox[\ifbool{latexml}{10em}{8em}][l]{\bfseries#1}}%
  \thmnumber{~#2}%
   \thmnote{\parbox[t]{.75\textwidth}{\bfseries\raggedright#3}}%
  }

\newtheoremstyle{apex}% name
  {0pt}% Space above, empty = `usual value'
  {0pt}% Space below
  {}% Body font
  {}% Indent amount (empty = no indent, \parindent = para indent)
  {\bfseries}% Thm head font
  {}% Punctuation after thm head
  {\newline}% Space after thm head
%  {\thmname{{\bfseries#1}}\thmnumber{{ \bfseries#2}}%
%   \thmnote{{\raggedright\bfseries#3}}%
  {\thmname{\makebox[\ifbool{latexml}{10em}{8em}][l]{\bfseries#1}}%
  \thmnumber{~#2}%
   \thmnote{\parbox[t]{\ifbool{latexml}{.6\textwidth}{.7\textwidth}}{\bfseries\raggedright#3}}%
  }% Thm head spec
  % the padding for the box takes just a bit of room away from these that examples get to keep

\theoremstyle{apexExample}
\newtheorem{example}{Example}[section]
\newcommand{\exampleautorefname}{Ex\-am\-ple}
\theoremstyle{apex}

%\makeatletter
%\renewenvironment{proof}[1][\proofname]{\pagebreak[2]\par
%  \pushQED{\qed}%
%  \normalfont \topsep6\p@\@plus6\p@\relax
%  \trivlist
%  \item[\hskip\labelsep
%        \bfseries
%    #1]\mbox{}\\* % something is needed to be able to get a newline
%}{%
%  \popQED\endtrivlist\@endpefalse
%}
%\makeatother
% this was causing: The number of automatic begin and end text-unit para hooks differ!

%\renewcommand{\qedsymbol}{\ensuremath{\square}}
% this was causing
%%%%Text line contains an invalid character.
%%%%l.10  <mi mathvariant="normal">^^03
%%%%                                   </mi>
% in the luamml-mathml.html file

\newspecialbox{definition}{Def\-i\-ni\-tion}{60}
% draw = yellow!95!black!60 = Hsb( 60,.59,.97)
% topc = white!95!yellow    = Hsb( 60,.05,1)
% botc = yellow!90!black!30 = Hsb( 60,.28,.97)

\newspecialbox{theorem}{The\-o\-rem}{120}
% draw = green!30!black!50  = Hsb(120,.23,.65)
% topc = white!95!green     = Hsb(120,.05,1)
% botc = green!60!black!20  = Hsb(120,.13,.92)

\newspecialbox{keyidea}{Key Idea}{0}
% draw = red!30!black!50    = Hsb(  0,.23,.65)
% topc = white!95!red       = Hsb(  0,.05,1)
% botc = red!60!black!20    = Hsb(  0,.13,.92)



\newtoggle{abridgeConics}
\toggletrue{abridgeConics}

\newcommand{\monthYear}{%
\ifcase \month \or January\or February\or March\or April\or May\or June\or July\or August\or September\or October\or November\or December\fi \space \number \year}
%modified from \today. we could do
%\usepackage[en-US]{datetime2}
%\DTMlangsetup{showdayofmonth=false}
% so that \today is just month and year
%but LaTeXML doesn't have datetime2, so we need this anyway

\usepackage{multirow}
%\pgfplotsset{width=\marginparwidth+1pt,compat=1.3}

%\usepackage{wrapfig}

\usepackage{booktabs}

\newlength{\cellwidth}

\setcounter{secnumdepth}{1}
\setcounter{tocdepth}{1}

\makeatletter
\let\ps@oldplain=\ps@plain % save the plain pagestyle
\makeatother

\usepackage{fancyhdr}

\renewcommand{\chaptermark}[1]{\markboth{\chaptername\ \thechapter\ \ \ \ {#1}}{}}
\renewcommand{\sectionmark}[1]{\markright{\thesection\ \ \ \  #1}}
\renewcommand{\headrulewidth}{0pt}
\renewcommand{\footrulewidth}{0pt}


\fancypagestyle{prose}{%
 \fancyhf{}
 \fancyhead[LE]{\nouppercase{\leftmark}}%
 \fancyhead[RO]{\nouppercase{\rightmark}}%
 \fancyfoot[LE]{\begin{minipage}{\textwidth}%
  \noindent\hspace{\marginparwidth}\hspace{\marginparsep}\hspace{-.4em}%
  \makebox[0pt][l]{\rule{\textwidth}{.4pt}}%
  \vskip.2\baselineskip%
  \noindent\hspace{\marginparwidth}\hspace{\marginparsep}\hspace{-.4em}%
  Notes:%
  \vskip 1.5in\textbf{\thepage}%
 \end{minipage}}

 \fancyfoot[RO]{\begin{minipage}{\textwidth+\marginparwidth+\marginparsep}%
  \rule{\textwidth-\marginparwidth-\marginparsep}{.4pt}
  \vskip.2\baselineskip
  Notes:
  \vskip 1.5in
  \hfill\textbf{\thepage}
 \end{minipage}}
 \fancyhfoffset[LE,RO]{\marginparsep+\marginparwidth}
}
\fancypagestyle{exercise}{%
	\fancyhf{}% 
	\fancyhfoffset[LE,RO]{32pt}%
	\fancyfoot[LE,RO]{\textbf{\thepage}}
}



\let\oldmainmatter\mainmatter
\renewcommand{\mainmatter}{%
 \oldmainmatter
 \fancypagestyle{plain}{% override the default for opening chapters
  \fancyhf{}
  \fancyfoot[RO]{\begin{minipage}{\textwidth+\marginparwidth+\marginparsep}%
   \rule{\textwidth-\marginparwidth-\marginparsep}{.4pt}
   \vskip.2\baselineskip
   Notes:
   \vskip 1.5in
   \hfill\textbf{\thepage}
  \end{minipage}}
  \fancyhfoffset[RO]{\marginparsep+\marginparwidth}
 }
 \pagestyle{prose}
}

\newtoggle{inappendix}
\makeatletter
\AddToHook{cmd/appendix/before}{%
 \let\ps@plain=\ps@oldplain% restore the pagestyle
 \cleardoublepage
}
\makeatother
\AddToHook{cmd/appendix/after}{%
 \toggletrue{inappendix}
 \setcounter{secnumdepth}{-1}
 \pagenumbering{arabic}
 \renewcommand{\thepage}{A.\arabic{page}}
 \renewcommand{\thechapter}{\arabic{chapter}}
 \part*{\appendixname}
}


% an enumerate like environment that can be mixed into tabular, array, etc.
\newcounter{anywhereenumi}
\newenvironment{anywhereenum}{%
 \setcounter{anywhereenumi}{0}%
 \renewcommand{\item}[1][]{%
  \ifx.##1.%
  \refstepcounter{anywhereenumi}%
  \makebox[1em][r]{\arabic{anywhereenumi}.}~~%
  \else%
  \makebox[1em][r]{##1.}~~%
  \fi%
 }%
}{}

\newcommand{\ds}{\displaystyle}

\newcommand{\primeskip}{\ifbool{mmode}{\mkern1.35mu}{\kern.075em}\relax}
%\newcommand{\primeskip}{\hskip.75pt}

\newcommand{\fp}{\ensuremath{f\,'}}
\newcommand{\fpp}{\ensuremath{f\,''}}

\newcommand{\Fp}{\ensuremath{F\primeskip'}}
\newcommand{\Fpp}{\ensuremath{F\primeskip''}}

\newcommand{\yp}{\ensuremath{y\primeskip'}}
\newcommand{\gp}{\ensuremath{g\primeskip'}}

\newcommand{\dd}{\operatorname{d}\!}

\newcommand*{\abs}[1]{\ensuremath{\left\lvert #1 \right\rvert}}
\newcommand*{\norm}[1]{\ensuremath{\left\lVert #1 \right\rVert}}
\newcommand*{\vnorm}[1]{\ensuremath{\norm{\vec #1}}}
\newcommand{\bracket}[1]{\left\langle #1\right\rangle}
\newcommand*{\proj}[2]{\ensuremath{\text{proj}_{\,\vec #2}{\,\vec #1}}}

\newcommand{\vecE}{\ensuremath{\vec E}}
\newcommand{\vecF}{\ensuremath{\vec F}}
\newcommand{\vecG}{\ensuremath{\vec G}}
\newcommand{\vecT}{\ensuremath{\vec T}}
\newcommand{\vece}{\ensuremath{\vec e}}
\newcommand{\vecf}{\ensuremath{\vec f}}
\newcommand{\vecg}{\ensuremath{\vec g}}
\newcommand{\veci}{\ensuremath{\vec\imath}}
\newcommand{\vecj}{\ensuremath{\vec\jmath}}
\newcommand{\veck}{\ensuremath{\vec k}}
\newcommand{\vecl}{\ensuremath{\vec l}}
\newcommand{\vecn}{\ensuremath{\vec n}}
\newcommand{\vecr}{\ensuremath{\vec r}}
\newcommand{\vecu}{\ensuremath{\vec u}}
\newcommand{\vecv}{\ensuremath{\vec v}}
\newcommand{\vecw}{\ensuremath{\vec w}}
\newcommand{\vecx}{\ensuremath{\vec x}}
\newcommand{\vecy}{\ensuremath{\vec y}}
\newcommand{\vrp}{\ensuremath{\vec r\hskip1.25pt '}}
\newcommand{\vsp}{\ensuremath{\vec s\primeskip '}}
\newcommand{\vrt}{\ensuremath{\vec r(t)}}
\newcommand{\vst}{\ensuremath{\vec s(t)}}
\newcommand{\vvt}{\ensuremath{\vec v(t)}}
\newcommand{\vat}{\ensuremath{\vec a(t)}}

\newcommand{\grad}{\vec\nabla}

\newcommand{\underlinespace}{\underline{\phantom{xxxxxx}}}

\newcommand{\zerooverzero}{\dfrac{\makebox[0pt]{\text{`` }0\text{ ''}}}0\ \ }


\DeclareMathOperator{\sech}{sech}
\DeclareMathOperator{\csch}{csch}
\DeclareMathOperator{\Div}{div}
%\DeclareMathOperator{\grad}{grad}
\DeclareMathOperator{\curl}{curl}
\DeclareMathOperator{\divv}{div}

%\newcommand*{\sword}[1]{\textbf{#1}}

\newcommand{\LHequals}{\mathrel{\overset{\text{by LHR}}{=}}}

\newcommand{\surfaceS}{\ensuremath{\mathcal{S}}}


%\newspecialbox[notempty]{exvideo}{ignored}{240}
\AtBeginDocument{\makeStyles{exvideo}{240}}
\newcommand{\setexvideocolor}{%
 \definecolor{topexvideo}{Hsb}{240,.05,1}% %= hsl(#4,100,97.5)
 \definecolor{borderexvideo}{Hsb}{240,.3,1}% %= hsl(#4,90.5,68.4)
 \definecolor{bottomexvideo}{Hsb}{240,.15,1}% %= hsl(#4,81.9,83.4)
}
\newcommand{\setexvideobw}{%
 \definecolor{topexvideo}{Hsb}{0,0,1}% white
 \definecolor{bottomexvideo}{Hsb}{0,0,1}% white
 \definecolor{borderexvideo}{Hsb}{0,1,0}% black
}
\newcommand{\exvideo}[1]{%
 \begin{tcolorbox}[
   colframe=borderexvideo,
   beforeafter skip=\boxskipamount,
   interior style={top color=topexvideo, bottom color=bottomexvideo},
   sharp corners=all,
   notitle,
   width=\textwidth,
   enhanced,
   tcbox width=forced left
  ]%
  #1%
 \end{tcolorbox}%
}


% \jmtVideo{youtube code}{jmt url suffix}{actual title}
%\newcommand{\jmtVideo}[3]{\genVideo{#1}{http://patrickjmt.com/#2/}{#3}}

%\newcommand{\khanVideo}[3]{\genVideo[?utm_campaign=embed]{#1}{https://www.khanacademy.org/video/#2}{#3}}


\makeatletter
% \mfigure[graphicsoptions]{marginaliaoptions}{caption}{label}{file}
\newcommand{\mfigure}[5][]{%
	\mnote@pex{\iffalse}{#2}{%
		\centering\myincludegraphics[#1]{#5}%
		%\captionsetup{type=figure}
		\caption{#3}\label{#4}%
	}%
}

% \mtable[marginaliaoptions]{caption}{label}{contents}
\newcommand{\mtable}[4][]{%
	\mnote@pex{\iffalse}{#1}{%
		\centering\small#4%\captionsetup{type=figure}%
		\caption{#2}\label{#3}%
	}%
}

% \mnote[marginaliaoptions]{contents}
\newcommand*{\mnote}[2][]{%
	\mnote@pex{\iftrue}{#1}{#2}%
}

% testmargin[top|bottom] was interfering with the tagging, as was the mbox
\newcommand*{\mnote@pex}[3]{%
	\marginalia[ysep page bottom=.8in,#2]{%[yshift=#2]
%	\marginnote{%
%		\begin{minipage}[t]{\marginparwidth}
			\captionsetup{type=figure}%
%			\testmargintop%
%			\mbox{}%\makebox[0pt]{\tikz\draw(0,0)circle(1pt);}%
%			\\[-\baselineskip]%
			#1\RaggedRight\fi#3\ifhmode\unskip\fi%
%			\testmarginbottom%\makebox[0pt]{\tikz\draw(0,0)circle(2pt);}
%		\end{minipage}%
%	}[#2]%
	}%
}
\makeatother


%\newenvironment{lxfigure}{%
%	\iflatexml%
%		\begin{figure}[!h]%
%	\else%
%		\noindent\begin{minipage}[t]{\linewidth}\noindent%
%	\fi%
%	\captionsetup{type=figure}%
%}{%
%	\iflatexml\end{figure}\else\end{minipage}\fi%
%}

\newcommand{\tbox}[1]{{\tagpdfsetup{table/tagging=presentation}\begin{tabular}{c}#1\end{tabular}}} % a tall box
\newcommand*{\zbox}[1]{\makebox[0pt][c]{#1}} % a zero width box






\newtoggle{isEarlyTrans}
\togglefalse{isEarlyTrans}

\newcommand{\prereqIntro}{The material in this section provides a basic review of and practice problems for pre-calculus skills essential to your success in Calculus. You should take time to review this section and work the suggested problems (checking your answers against those in the back of the book). Since this content is a pre-requisite for Calculus, reviewing and mastering these skills are considered your responsibility. This means that minimal, and in some cases no, class time will be devoted to this section. When you identify areas that you need help with we strongly urge you to seek assistance outside of class from your instructor or other student tutoring service.\bigskip}

\newcommand*{\inputprereq}[1]{%
  \addtocounter{section}{-1}%
  \addtocontents{toc}{\protect\setcounter{tocdepth}{0}}%
  \input{#1}%
  \addtocontents{toc}{\protect\setcounter{tocdepth}{1}}%
}

% lets try to reduce bad boxes
\usepackage{microtype}
\hfuzz=2pt\relax
\vfuzz=1.5\baselineskip
% ignore overfull < this amount
%\newdimen\hfuzz % lock it in?
%\newdimen\vfuzz % lock it in?
%\hbadness=10000
\vbadness=9999\relax
% ignore underfull > this amount
\parskip=0pt plus \baselineskip
\baselineskip=1\baselineskip plus .3\baselineskip

\ifbool{xetex}%
	{%
	\sffamily
%	\usepackage{fontspec}
%	\usepackage{unicode-math}
	\usepackage{mathspec}
	\setallmainfonts[Mapping=tex-text]{Calibri}
	\setmainfont[Mapping=tex-text]{Calibri}
	% setallmainfonts claims to setmainfont. but it doesn't?
%	\setmathsfont[Mapping=tex-text]{Calibri}
%	\setmathrm[Mapping=tex-text]{Calibri}
	\setsansfont[Mapping=tex-text]{Calibri}
	\setmathsfont(Greek){[cmmi10]}
	}
	{}

\ifbool{luatex}%
	{%
	\sffamily
	\usepackage{fontspec}
	\usepackage{unicode-math}
	%\usepackage{mathspec}
	%\setallmainfonts[Mapping=tex-text]{Calibri}
%	\setmainfont{Calibri}
	%\setsansfont[Mapping=tex-text]{Calibri}
%	\setmathfont{latinmodern-math.otf}
	% Calibri doesn't quite seem set for math
%	\setmathfont[range=\mathup]{Calibri}
%	\setmathfont[range=\mathit]{Calibri Italic} % missing characters
%	\usepackage[default]{lato} % Lato/lete does not give a valid unicode vector arrow
	\usepackage{lete-sans-math}
	\usepackage[default,defaultsans]{lato}
	\usepackage{realscripts}
	% Todo bug See https://tex.stackexchange.com/a/748076/107497
	\DeclareMicrotypeAlias{LeteSansMath.otf}{TU-basic}
	\let\pdfsavepos=\savepos
	\let\pdflastxpos=\lastxpos
	\let\pdflastypos=\lastypos
	}
	{}

\ifbool{latexml}{
 \usepackage[american]{babel}
}{
 \usepackage{polyglossia}
 \setdefaultlanguage[variant=usmax]{english}
 \renewcommand*{\englishhyphenmins}{22}
 \AfterEndPreamble{
 \hyphenation{%
  an-ti-der-iv-a-tive
  an-ti-der-iv-a-tives
  app-rox-i-mate
  cen-tered
  chang-es
  con-struc-tions
  de-creas-es
  Der-iv-a-tive
  der-iv-a-tive
  dis-place-ment
  dis-tance
  e-qual-ly
  ex-am-ples
  Func-tions
 % Hô-pi-tal % doesn't hyphenate L'Hôpital's
  im-pli-cit
  in-dis-tin-guish-a-ble
  in-fall-i-ble
 % %L'Hô-pi-tal's % ' causes: ! Not a letter.
 % % see https://tex.stackexchange.com/a/165091/107497 for fix and pitfalls
  meth-od
  of-ten
  proc-ess
  re-fer-ring
  qua-dra-tic
  sa-li-ent
  se-quence
  sketch-ing
  smart-er
  sub-sti-tute
  The-o-rem
  Trig-o-no-me-tric
  trig-o-no-me-tric
  wheth-er
 }}
}


\usepackage{tocbibind}
%\let\oldprintindex\printindex
%\renewcommand{\printindex}{%
% \cleardoublepage
%% \chapter{\indexname} % \printindex has its own heading
% \phantomsection
%% \iflatexml\chapter*{\indexname}\fi
% \addcontentsline{toc}{chapter}{\indexname}
% \oldprintindex
%}


\newtoggle{bsc} % default false
\newcommand{\forwhom}{\iftoggle{bsc}{ for Bismarck State College}{}}

\title{APEX Calculus LT}
\author{Greg Hartman, VMI \and UND Math Dept}
\usepackage[
	bookmarksnumbered,
	hidelinks,
	pdfstartview=FitH,
	linktoc=all,
	pdfdisplaydoctitle,
	bookmarksdepth=2,
]{hyperref}
\hypersetup{
	pdftitle={APEX Calculus LT},
	pdfauthor={UND Math Dept and Greg Hartman, VMI},
	unicode,
    pdflang=en-US
}
%\usepackage{hyperxmp}
\ifbool{latexml}{}{
 \usepackage{bookmark}
}


% hyperref changes these
% if they come before and have newcommand, latexml overwrites them
\AtBeginDocument{
 \renewcommand{\chapterautorefname}{Chap\-ter} % the default is lowercase
 \renewcommand{\sectionautorefname}{Sec\-tion} % the default is lowercase
 \renewcommand{\figureautorefname}{Fig\-ure}
 \renewcommand{\appendixname}{Ap\-pen\-di\-ces}
}
\newcommand{\exampleEnvautorefname}{Ex\-am\-ple}
\newcommand{\autoeqref}[1]{\hyperref[#1]{\equationautorefname~(\ref*{#1})}}
% autoref doesn't use parentheses

% \apex has to be *used* after hyperref
% lxNavbar has to come after latexml
\begin{lxNavbar}
\lxRef{lxApexTOC}{Table of Contents}\\
\lxContextTOC
\end{lxNavbar}

\lxRequireResource{LaTeXML-marginpar.css}
\lxRequireResource{LaTeXML-navbar-left.css}
\lxRequireResource{style.css}
\lxRequireResource{style-theorems.css}
\lxRequireResource[media=(max-width:1000pt)]{style-narrow.css}
\lxRequireResource{%
https://ajax.googleapis.com/ajax/libs/jquery/1.12.2/jquery.min.js}
\lxRequireResource{script.js}
\lxRequireResource{LaTeXML-maybeMathJax4.js}

% set the defaults, just in case
\printincolor
\usetwoDgraphics


% todo bugfix see https://github.com/latex3/tagging-project/issues/981
%\newcommand { \intertextmathbegin } { \pdf_version_compare:NnF<{2.0} {
% \tag_struct_begin:n { tag=Formula \bool_if:NT \l__tag_math_alt_bool { , alt={subformula~of~larger~formula} } }
%} }
%\newcommand { \intertextmathbegin } { \pdf_version_compare:NnF<{2.0} { \tag_struct_begin:n{tag=Formula} } }
%\newcommand { \intertextmathend } { \pdf_version_compare:NnF<{2.0} { \tag_struct_end: } }
%\newcommand{\intertextmath}[1]{\intertextmathbegin\(#1\)\intertextmathend}
\newcommand{\intertextmath}[1]{\(#1\)}
% see https://github.com/latex3/tagging-project/issues/1396
%\newcommand{\intertextmath}[1]{}

\makeatletter
\newcommand{\exercisesubsection}[1]{%
 \clearpage
 \@startsection{subsection}{2}{-1sp}{0pt}{\baselineskip}{\huge\bfseries}*{\texorpdfstring{\hyperref[sol#1]{Exercises}}{Exercises}}
 \hrule
 \noindent
}
%\newcommand{\exercisesubsection}[1]{%
% \bgroup
%  \titleformat{\subsection}{\huge\bfseries}{}{0pt}{}[\hrule\vspace{-\baselineskip}]%
%  \subsection{\texorpdfstring{\hyperref[sol#1]{Exercises}}{Exercises}}%
% \egroup
%}

% this seems higher level, but tagpdf isn't liking the spare /Refs running around
%\titleclass{\exercisesubsection}{straight}[\section]
%\titleformat{\exercisesubsection}[block]{\clearpage\setcounter{enumXi}{0}}%format
% {}%label
% {0pt}%sep
% {\exercisesubsectiontitle}%before
% %after
%\titlespacing*{\exercisesubsection}{0pt}{0pt}{-7ex}
%\dottedcontents{exercisesubsection}[6.1em]{}{3.2em}{1pc}
% exercisesubsection doesn't appear in the toc (so I don't care about the numbers), but without this the subsequent section gets messed up
% at is still letter
\newcommand{\exercisesubsubsection}[1]{%
 \noindent
 \@startsection{subsubsection}{3}{0pt}{-.5\bigskipamount}{\baselineskip}{\Large\itshape}*{#1}%
% \par
}
%\newcommand{\exercisesubsubsection}[1]{%
% \bgroup
%  \titleformat{\subsubsection}[block]{\noindent}%format
%   {is not labeled}%label
%   {0pt}%sep
%   {\noindent\Large\textit}%before
%   %after
%   \titlespacing*{\subsubsection}{0pt}{.5\bigskipamount}{0pt}
%  \subsubsection*{#1}
% \egroup
%}
%\titleclass{\exercisesubsubsection}{straight}[\exercisesubsection]
%\titleformat{\exercisesubsubsection}[block]{\noindent}%format
% {is not labeled}%label
% {0pt}%sep
% {\noindent\Large\textit}%before
% %after
%\titlespacing*{\exercisesubsubsection}{0pt}{.5\bigskipamount}{0pt}
%\dottedcontents{exercisesubsubsection}[6.1em]{}{3.2em}{1pc}
% exercisesubsubsection doesn't appear in the toc (so I don't care about the numbers), end of story
\makeatother


\InputIfFileExists{./options}{}{%
 \newcommand{\thetitle}{Calculus}
 \printallanswers
 \printincolor
% \usetwoDgraphics
 \usethreeDgraphics
}

% these allow us to skip pdf tagging for
% https://tex.stackexchange.com/a/751610
\ExplSyntaxOn
\providecommand{\tag_struct_begin:n}[1]{}
\providecommand{\tagstructbegin}[1]{}
\providecommand{\tag_struct_parent_int:}{}
\providecommand{\pdfannot_box_ref_last:}{}
\providecommand{\tag_struct_insert_annot:nn}[2]{}
\providecommand{\tag_struct_end:}{}
\providecommand{\tagstructend}{}
\providecommand{\tag_mc_begin:n}[1]{}
\providecommand{\tagmcbegin}[1]{}
\providecommand{\tag_mc_end:}{}
\providecommand{\tagmcend}{}
\AtBeginDocument{\providecommand{\luamml_ignore:}{}}
\ExplSyntaxOff
\providecommand{\tagpdfsetup}[1]{}
\providecommand{\ResumeTagging}[1]{}
\providecommand{\SuspendTagging}[1]{}
\providecommand{\tagpdfparaOn}{}
\providecommand{\tagpdfparaOff}{}
\providecommand{\tagtool}[1]{}

\pgfkeysifdefined{/tikz/alt}{}{
 \tikzset{alt/.code={}}
 \tikzset{actualtext/.code={}}
}
